% NOTE: Results below are from synthetic data and MUST be replaced with live evaluation results.
% See PUBLICATION_GUIDE.md for complete instructions on generating real results.
% DO NOT SUBMIT until live evaluation is complete and all numbers are updated.
\documentclass[11pt,a4paper]{article}
\usepackage[utf8]{inputenc}
\usepackage{amsmath,amssymb,amsfonts}
\usepackage{booktabs}
\usepackage{graphicx}
\usepackage{hyperref}
\usepackage[margin=1in]{geometry}
\usepackage{cite}

\hypersetup{
    colorlinks=true,
    linkcolor=blue,
    citecolor=blue,
    urlcolor=blue
}

\title{\textbf{Beyond Retrieval: Evaluating Evidence-Grounded Agentic AI Systems for Reliable Technical Question Answering}}

\author{
    \textbf{Gyanankur Pratim Baruah} \\
    \textit{B.Sc. Data Science, Third Year (TY)} \\
    Thakur Shyamnarayan Degree College, Mumbai University \\
    \textit{Student Research Fellow \& Lead Author} \\
    \texttt{gyanankur9@gmail.com}
    \and
    \textbf{Anjali Tiwari} \\
    \textit{B.Sc. Data Science, Third Year (TY)} \\
    Thakur Shyamnarayan Degree College, Mumbai University \\
    \textit{Student Research Fellow \& Co-Author} \\
    \texttt{anjalitiwari30082006@gmail.com}
}

\date{\today}

\begin{document}

\maketitle

\begin{abstract}
Retrieval-Augmented Generation (RAG) pipelines often fail when deployed in technical domains characterised by temporal version drift, contradictory documentation, and incomplete corpora. Conventional single-pass retrieval architectures frequently reproduce deprecated API syntax and hallucinate under-specified answers. To address these limitations, we introduce an evidence-grounded, multi-agent architecture incorporating version-aware metadata routing, claim-level verification, contradiction detection, and selective abstention. Evaluated across an independent 100-question technical benchmark stratified across five adversarial conditions (Python, FastAPI, LangChain, ChromaDB, Pandas, and GitHub), our proposed system achieves a statistically significant improvement in answer accuracy ($88.0\%$, $p < 0.001$) and citation precision ($0.877$) compared to conventional baselines. Our ablation findings demonstrate that explicit contradiction handling and calibrated abstention are indispensable for dependable technical question answering, though multi-agent orchestration introduces predictable trade-offs in latency and token expenditure.
\end{abstract}

\section{Introduction}
Large language models (LLMs) have transformed software engineering workflows, yet their reliance on parametric memory leads to hallucinations and factual obsolescence when querying rapidly evolving technical documentation. Retrieval-Augmented Generation (RAG) mitigates this by grounding generation in external document stores. However, standard RAG pipelines—relying on single-pass vector similarity search followed by direct generation—exhibit severe reliability vulnerabilities when facing real-world software engineering challenges:
\begin{enumerate}
    \item \textbf{Version Drift:} Outdated legacy APIs mixed with current specifications in vector indices lead models to reproduce deprecated syntax.
    \item \textbf{Source Contradiction:} Conflicting guidance across tutorials, issues, and official documentation cause unhandled contradictions.
    \item \textbf{Corpus Incompleteness:} When retrieved evidence is insufficient, standard models attempt speculative extrapolation rather than graceful refusal.
\end{enumerate}

To bridge this gap, this paper investigates whether specialised multi-agent verification and evidence-management controls improve technical question answering beyond conventional RAG. We introduce a controlled experimental framework comparing four architectural configurations over a rigorous 100-question technical benchmark.

\section{Related Work}
Recent scholarship has increasingly focused on the limitations of naive RAG architectures. Chen et al.~\cite{chen2025multiagent} examined multi-agent collaboration failure modes and communication bottlenecks, noting high error propagation rates in unverified agent chains. Zhao et al.~\cite{zhaoo2025versionrag} introduced VersionRAG and the VersionQA benchmark to address outdated document retrieval, though their scope is limited to single-agent version routing without multi-agent claim verification. Concurrently, Patel et al.~\cite{patel2026techdocrag} explored evidence tracing in technical documentation, and Smith et al.~\cite{smith2026self} studied iterative self-correction. Finally, Kumar et al.~\cite{kumar2026uncertainty} investigated uncertainty calibration and refusal behavior in RAG, though without multi-agent orchestration or contradiction detection. Recent work on contradiction detection~\cite{contradiction2025rag,conflicts2025rag} and conformal abstention~\cite{conformal2024abstention,knowlimits2025survey} provides complementary approaches to managing conflicting evidence and selective refusal.

\section{Research Gap}
While prior works examine isolated aspects of version routing~\cite{zhaoo2025versionrag} or multi-agent collaboration failure~\cite{chen2025multiagent}, existing literature lacks a unified, evidence-grounded framework that simultaneously evaluates:
\begin{itemize}
    \item Claim-level provenance attribution across dependent technical guides.
    \item Automated temporal version conflict resolution.
    \item Calibrated selective abstention under insufficient evidence.
    \item Controlled ablation of multi-agent overhead versus epistemic reliability gains.
\end{itemize}
This paper directly addresses these gaps.

\section{Methodology}
We evaluate four system configurations over 100 benchmark questions: \textbf{E1} (Conventional RAG), \textbf{E2} (RAG + Single Verifier), \textbf{E3} (Multi-Agent RAG), and \textbf{E4} (Proposed Evidence-Aware Multi-Agent System). 

Our benchmark (`data/questions.json`) is evenly stratified across five conditions ($n=20$ each): direct factual retrieval, multi-document synthesis, outdated document resolution, contradictory sources, and insufficient evidence (unanswerable queries).

Configuration E4 orchestrates six specialised agents (Query Analyst, Retrieval Agent, Evidence Analyst, Conflict Detector, Verifier, and Response Controller) within a LangGraph state machine to enforce rigorous grounding and selective abstention ($\text{ABSTAIN:}$).

\section{Implementation}
The framework is implemented in Python using LangGraph (v0.2+), ChromaDB (v0.4+), and Sentence Transformers (`all-MiniLM-L6-v2`). Technical documents are chunked (500 tokens, 50 overlap) with explicit version and source metadata. Evaluation harnesses record token consumption, latency, and correctness.

\section{Results}
Table~\ref{tab:results} summarizes overall performance across configurations ($N=100$).

\begin{table}[htbp]
\centering
\caption{System Performance Across Configurations (Mean $\pm$ 95\% CI)}
\label{tab:results}
\begin{tabular}{lccccc}
\toprule
\textbf{Configuration} & \textbf{Accuracy (\%)} & \textbf{Citation Prec.} & \textbf{Latency (s)} & \textbf{Cost / Query (\$)} \\
\midrule
E1: Conventional RAG & $68.0 \pm 9.3$ & $0.518 \pm 0.043$ & $1.20 \pm 0.04$ & $\$0.0020$ \\
E2: Single Verifier & $78.0 \pm 8.3$ & $0.655 \pm 0.036$ & $2.10 \pm 0.06$ & $\$0.0050$ \\
E3: Multi-Agent RAG & $79.0 \pm 8.1$ & $0.708 \pm 0.037$ & $3.79 \pm 0.10$ & $\$0.0120$ \\
E4: Proposed Evidence-Aware & \textbf{88.0 $\pm$ 6.5} & \textbf{0.877 $\pm$ 0.026} & $4.47 \pm 0.13$ & $\$0.0179$ \\
\bottomrule
\end{tabular}
\end{table}

Paired t-tests confirm that Configuration E4 outperforms E1 significantly in accuracy ($t = 3.76, p < 0.001$) and citation precision ($t = 14.24, p < 0.001$). Ablation analysis reveals that removing the Selective Abstention Controller inflicts the largest accuracy drop ($-14.0\%$).

\section{Discussion}
The results demonstrate a clear trade-off between epistemic reliability and computational overhead. While E4 improves accuracy to $88.0\%$, it incurs a $3.72\times$ increase in latency and $8.95\times$ token expenditure. Production systems should employ hybrid routing to invoke multi-agent verification selectively.

\section{Limitations}
Limitations include corpus restriction to open-source software documentation, underlying model nondeterminism, and cumulative multi-agent execution latency.

\section{Conclusion}
We introduced an evidence-grounded multi-agent RAG architecture that resolves version drift, contradictions, and insufficient evidence. Our findings establish that explicit contradiction handling and calibrated abstention are vital for dependable technical question answering.

\section*{Acknowledgments}
The authors gratefully acknowledge the institutional support and computing resources provided by the Department of Data Science at Thakur Shyamnarayan Degree College, Mumbai University. We also thank our faculty mentors and student research peers for their constructive feedback during the benchmark annotation and experimental validation phases.

\bibliographystyle{plain}
\bibliography{references}

\end{document}
